\section{Strong recentering and the full current spectrum}\label{sec:recentering}
\subsection{A retained source with ordinary \Ltwo{} pairing}
For the complete vector field in \eqref{eq:current},
\begin{equation}
Y\theta=-\sin\theta,\qquad
\tan(\theta_t/2)=e^{-t}\tan(\theta/2),\qquad
\operatorname{div}_mY=-3\cos\theta.
\end{equation}
These follow from the four distinct link derivatives of P.
Consequently
\begin{equation}\label{eq:flow-jacobian}
j_t(U)=\left(\frac{\sin\theta_t}{\sin\theta}\right)^3
                     \frac{w(\Phi_tU)}{w(U)} .
\end{equation}
For $P\ne-I$ the trajectory converges as $t\to+\infty$:
its remaining length is $\theta_t/\sqrt\ell$, since
$\norm Y=\sin\theta/\sqrt\ell$.
Write $\Phi_\infty U$ for this limit and
$w_\infty(U)=w(\Phi_\infty U)$. This is a bounded positive
measurable function. Smoothness at the repelling critical set
is not asserted.

\begin{theorem}[Native strong recovery]\label{thm:recentering}
The actual vectors $J_Rf=\UU_RS_Rf$ converge strongly to
\begin{align}
J_{\rm rec}f(U)
 &=\sqrt{\frac{w_\infty(U)}{\rho(0)w(U)}}\,
                    J_{\rm H}f(\theta(U))\,I,\label{eq:rec-source}\\
J_{\rm H}f(\theta)
 &=\frac{(1+y^2/4)^{3/2}}{y}
           \int_0^\infty H_f(t)\sin(yt)\dd t,
\qquad y=2\tan(\theta/2).
\end{align}
Moreover
\begin{equation}\label{eq:rec-pair}
\ip{J_{\rm rec}f}{J_{\rm rec}g}_\nu=\ip f g_2,\qquad
J_{\rm rec}U_t=\UU_{-t}J_{\rm rec}.
\end{equation}
Thus $J_{\rm rec}$ extends to an isometry from $L^2(\R)$ into
the same physical boundary completion.
\end{theorem}
\begin{proof}
First use Haar. The map
\begin{equation}
(\mathcal RF)(y)=\sqrt{2/\pi}\,
\frac{y}{(1+y^2/4)^{3/2}}F(2\arctan(y/2))
\end{equation}
is a unitary from radial Haar \Ltwo{} to $L^2(\R_+,\dd y)$.
Apply it to the Haar recentered packet and use
\eqref{eq:flow-jacobian}. The exact result is
\begin{equation}\label{eq:rec-riemann}
\frac{\sqrt{2/\pi}}{N\sqrt{1+(y/(2N))^2}}
\sum_{n\ge1}H_f(n/N)
          \sin\!\left(2n\arctan(y/(2N))\right).
\end{equation}
On bounded y intervals the Riemann sums converge uniformly to
$\sqrt{2/\pi}\int H_f(t)\sin(yt)\dd t$.
The squared norm before this limit is exactly
$N^{-1}\sum|H_f(n/N)|^2\to\norm{H_f}_2^2$.
The sine transform is unitary with this normalization, so
local convergence, boundedness, and equality of limiting norms
give strong convergence on the whole half-line.
Inverting $\mathcal R$ gives $J_{\rm H}$.

For the actual measure, $\UU_RS_Rf$ equals the Haar expression
times $\sqrt{w(\Phi_RU)/(\rho(0)w(U))}$. These multipliers
are uniformly bounded and converge almost everywhere to the
factor in \eqref{eq:rec-source}. Strong convergence follows
by dominated convergence and the Haar strong limit. Unitarity
of $\UU_R$ and \eqref{eq:phase-orthogonality} prove the pairing.
Finally $S_RU_tf=S_{R+t}f$ exactly. Passing to the strong
limit in
$\UU_RS_{R+t}f=\UU_{-t}(\UU_{R+t}S_{R+t}f)$ proves covariance.
\end{proof}
The formula describes actual full-state dependence through
$w_\infty/w$. It does not assume that the current preserves
the class sector. Since $\UU_R$ is unitary,
$\UU_RF_Rf$ has the same strong limit.

Winding retains only part of its norm under this recentering:
\begin{equation}\label{eq:rec-winding-weak}
\UU_RV_aS_Rf\rightharpoonup
a^{-1/2}J_{\rm rec}U_{\log a}f ,
\end{equation}
while the missing limiting squared norm is
\begin{equation}\label{eq:rec-defect}
\left(\frac{\rho_a(0)}{\rho(0)}-\frac1a\right)\norm f_2^2
=\frac{\sum_{j=1}^{a-1}\rho(2\pi j/a)}
       {a\rho(0)}\norm f_2^2>0\quad(a\ge2).
\end{equation}
Indeed each nonidentity branch, after the flow, tends to zero
on compact sets away from $P=-I$ by summation by parts in
its oscillatory coefficient series. Boundedness and density
give weak convergence to zero there. Its retained norm escapes
toward the repelling set, which has measure zero. Equations
\eqref{eq:branch} and \eqref{eq:wound-norm} then prove
\eqref{eq:rec-winding-weak}--\eqref{eq:rec-defect}.
This recentering solves weak escape for the unbranched \Ltwo{}
source, but does not retain the full winding module or the
Weil norm.

\subsection{An exact normal form for the entire current}
\begin{theorem}[Full interacting current is absolutely continuous]\label{thm:current-spectrum}
There is a transverse covariant Hilbert space $\cK$ and a unitary
identification
\begin{equation}\label{eq:current-normal}
\cH\simeq L^2(\R,\dd r;\cK),\qquad
\UU_tH(r)=H(r-t),\qquad A_\nu=\ii\partial_r ,
\end{equation}
with the usual Hilbert-valued $H^1$ generator domain.
In particular every finite-vector coefficient satisfies
\begin{equation}\label{eq:current-decay}
\ip\xi{\UU_t\eta}_\nu\longrightarrow0\quad(|t|\to\infty).
\end{equation}
\end{theorem}
\begin{proof}
Off $P=\pm I$, put $r=\log\tan(\theta/2)$, so $Yr=-1$.
Every remaining trajectory crosses the regular hypersurface
$\Sigma=\{\theta=\pi/2\}$ once. The map
$(r,z)\mapsto\Phi_{-r}z$ is a smooth bijection with smooth
inverse from $\R\times\Sigma$ onto this full-measure part
of M. The excluded sets have Haar measure zero by the
plaquette submersion, and therefore have $\nu$ measure zero.
This coordinate statement includes every link, not just the
plaquette angle.

Choose a smooth positive gauge-invariant transverse measure
$\sigma$, and write the pulled-back measure as
$b_\nu(r,z)\dd r\dd\sigma(z)$. Multiplication by
$b_\nu^{1/2}$ changes it to product measure. The weighted
flow factor is exactly
$\sqrt{b_\nu(r-t,z)/b_\nu(r,z)}$, so this multiplication
converts the group into translation. Gauge transformations
preserve r and commute with the flow, restricting only the
transverse Hilbert space. Fourier transformation in r
proves pure absolute continuity. Each matrix coefficient is
the Fourier transform of an $L^1$ function, by Hilbert-valued
Cauchy--Schwarz, and hence tends to zero.
\end{proof}
On the recovered subspace \eqref{eq:rec-pair} implies
$A_\nu J_{\rm rec}=J_{\rm rec}\Px$ on the smooth-test domain.
The full theorem is stronger than that subspace statement:
changing to an arbitrary seed vector elsewhere in the physical
boundary space cannot turn this current into a pure point
arithmetic translation generator.
